\section{时间反演流与对偶性}
\label{sec:time-reversal-flow-duality}

本节研究\personFloer{}复形中的两种对偶配对：一种由流的时间反演诱导，不使用环境流形的
定向；另一种由\personPoincare{}对偶诱导，本质上使用环境流形的定向。

\subsection{时间反演对偶}
\label{subsec:time-reversal-duality}

本小节通过明确考察对\personMorse{}函数及定向系统的依赖性，比较经典\personMorse{}复形上的两种
不同相容定向系统。这里不使用 $M$ 的可定向性。

设
\[
 \sigma(f)=\{o_\sigma(x;f);o_\sigma(x,y;f)\}_{x,y\in\operatorname{Crit}f}
\]
是第 15.6.1 节定义的 $f$ 的\personMorse{}复形上的相容定向系统：\emph{先定向不稳定流形，
再给不稳定流形配上余定向}。它诱导\personMorse{}边界映射
$\partial_f^\sigma:CM_k(f)\to CM_{k-1}(f)$。以
\[
 \delta_f:CM^*(f)\longrightarrow CM^{*+1}(f)
\]
表示对偶复形的相应上边界映射，其中
\[
 CM^\ell(f)=\operatorname{Hom}(CM_\ell(f),\mathbb Z).
\]

回忆，作为集合有 $\operatorname{Crit}f=\operatorname{Crit}(-f)$，而 $f$ 与 $-f$
的稳定、不稳定流形互换。因此，对每个
$y\in\operatorname{Crit}(-f)$，在
\[
 T_yW_f^u(y)=T_yW_{-f}^s(y)
\]
上赋予定向 $\sigma$。注意，作为集合有 $W_{-f}^s(y)=W_f^u(y)$。

接着，以 $\widetilde\sigma(g)$ 表示如下定向系统：\emph{先定向 $g$ 的稳定流形，再为不稳定
流形配上余定向}。

\begin{lemma}
\label{lem:time-reversed-orientation-coherent}
若 $\sigma(f)$ 相容，则定向系统
\[
 \widetilde\sigma(-f)
 =\{o_{\widetilde\sigma}(y;-f)\}_{y\in\operatorname{Crit}(-f)}
\]
相容。
\end{lemma}

\begin{exercise}
\label{ex:prove-time-reversed-orientation-coherent}
证明此引理。
\end{exercise}

利用这一系统，对每对 $y,x\in\operatorname{Crit}(-f)$，也在
$\mathcal M_g(y,x;-f)$ 上诱导一个定向，进而得到边界映射
\[
 \partial_{-f}^{\widetilde\sigma}:CM_\ell(-f)\longrightarrow CM_{\ell-1}(-f).
\]
按定义，“反演流”给出典范一一对应
\[
 \mathcal M_g(x,y;f)\longleftrightarrow\mathcal M_g(y,x;-f).
\]
特别地，$\dim\mathcal M_g(x,y;f)=0$ 当且仅当
$\dim\mathcal M_g(y,x;-f)=0$。此时可以比较两个整数
$n^\sigma(x,y;f)$ 与 $n^{\widetilde\sigma}(y,x;-f)$。

回忆有典范同构
\[
 CM^\ell(f)\longrightarrow CM_{n-\ell}(-f),
\]
它是基元素赋值 $x^*\mapsto\widetilde x$ 的线性延拓；这里
$x\in\operatorname{Crit}f$、$\mu_f(x)=\ell$，$x^*$ 是其对偶，而
$\widetilde x$ 是把 $x$ 看作 $-f$ 的临界点。

\begin{proposition}
\label{prop:time-reversal-morse-sign}
假设 $\mu_f(x)-\mu_f(y)=1$ 且 $\mu_f(y)=k$，则
\[
 n^{\widetilde\sigma}(y,x;-f)=(-1)^k n^\sigma(x,y;f).
\]
特别地，有交换图
\begin{equation}
\begin{array}{ccc}
 CM^k(f)&\xrightarrow{\;(-1)^k\delta_f\;}&CM^{k+1}(f)\\[4pt]
 \big\downarrow&&\big\downarrow\\[4pt]
 CM_{n-k}(-f)&\xrightarrow{\;\partial_{-f}^{\widetilde\sigma}\;}&
 CM_{n-k-1}(-f).
\end{array}
\label{eq:19-time-reversal-commutative-diagram}
\end{equation}
\end{proposition}

\begin{proof}
令 $p\in S W_{-f}^u(y)\cap S W_{-f}^s(x)$。点 $p$ 关于相容定向
$\widetilde\sigma$ 的符号对应于 $o_{\widetilde\sigma}(y,x;-f)$，并由定向向量空间等式
\begin{equation}
 T_p^{\widetilde\sigma}(S W_{-f}^u(y))
 =o_{\widetilde\sigma}(y,x;-f)\,
 N_p^{\widetilde\sigma}(S W_{-f}^s(x))
\label{eq:19-time-reversal-sign-definition}
\end{equation}
决定。这里参见式 (15.6.46) 附近的讨论。
\footnote{译者注：原书在本证明的定向符号中反复写
$o_{\widetilde\sigma}(x,y;-f)$，但相应反演流线是从 $y$ 到 $x$，命题陈述及前文计数
均写 $n^{\widetilde\sigma}(y,x;-f)$；译文统一为 $o_{\widetilde\sigma}(y,x;-f)$。}

另一方面，按 $o_\sigma(x,y;f)$ 的定义，
\begin{equation}
 N_p^\sigma(S W_f^s(y))
 =o_\sigma(x,y;f)\,T_p^\sigma(S W_f^u(x)).
\label{eq:19-original-flow-sign-definition}
\end{equation}
现在比较~\eqref{eq:19-time-reversal-sign-definition} 与
\eqref{eq:19-original-flow-sign-definition}。

式~\eqref{eq:19-time-reversal-sign-definition} 等价于
\begin{equation}
\begin{aligned}
 \{\nabla f(p)\}\oplus N_p^{\widetilde\sigma}(W_{-f}^s(x))
 &=o_{\widetilde\sigma}(y,x;-f)
 \{\nabla f(p)\}\oplus T_p^{\widetilde\sigma}(S W_{-f}^u(y))\\
 &=o_{\widetilde\sigma}(y,x;-f)
 T_p^{\widetilde\sigma}W_{-f}^u(y).
\end{aligned}
\label{eq:19-time-reversal-expanded-orientation}
\end{equation}

另一方面，有如下引理。

\begin{lemma}
\label{lem:stable-minus-f-unstable-f-orientation}
\emph{作为定向向量空间}，
\begin{equation}
 T_p^{\widetilde\sigma}W_{-f}^s(x)=T_p^\sigma W_f^u(x).
\label{eq:19-stable-minus-f-unstable-f}
\end{equation}
\end{lemma}

\begin{proof}
由 $W_{-f}^s(x)$ 的余定向定义，作为定向向量空间有
\[
 N_p^{\widetilde\sigma}W_{-f}^s(x)
 \oplus T_p^{\widetilde\sigma}W_{-f}^s(x)=T_pM.
\]
另一方面，由 $\widetilde\sigma$ 的定义，
\[
 N_p^\sigma W_f^u(x)=N_p^{\widetilde\sigma}W_{-f}^s(x),
\]
故
\[
 N_p^{\widetilde\sigma}W_{-f}^s(x)\oplus T_p^\sigma W_f^u(x)=T_pM.
\]
比较两式即得结论。
\end{proof}

在式~\eqref{eq:19-time-reversal-expanded-orientation} 右侧直和
$T_p^{\widetilde\sigma}(S W_{-f}^s(x))$，得到
\[
\begin{aligned}
 (-1)^{k+1}T_pM
 &=o_{\widetilde\sigma}(y,x;-f)
 T_p^{\widetilde\sigma}W_{-f}^u(y)
 \oplus T_p^{\widetilde\sigma}(S W_{-f}^s(x))\\
 &=o_{\widetilde\sigma}(y,x;-f)
 \{\nabla f(p)\}\oplus T_p^{\widetilde\sigma}(S W_{-f}^u(y))
 \oplus T_p^{\widetilde\sigma}(S W_{-f}^s(x))\\
 &=o_{\widetilde\sigma}(y,x;-f)(-1)^{n-k-1}
 T_p^{\widetilde\sigma}(S W_{-f}^u(y))
 \oplus T_p^{\widetilde\sigma}W_{-f}^s(x).
\end{aligned}
\]
这里左端使用恒等式
\[
\begin{aligned}
 &o_{\widetilde\sigma}(y,x;-f)
 T_p^{\widetilde\sigma}W_{-f}^u(y)
 \oplus T_p^{\widetilde\sigma}(S W_{-f}^s(x))\\
 &\quad=(-1)^{k+1}T_p^{\widetilde\sigma}W_{-f}^u(y)
 \oplus o_{\widetilde\sigma}(y,x;-f)
 \oplus T_p^{\widetilde\sigma}(S W_{-f}^s(x))\\
 &\quad=(-1)^{k+1}T_pM,
\end{aligned}
\]
右端也类似。

在式~\eqref{eq:19-original-flow-sign-definition} 右侧直和
$T_p^\sigma W_f^s(y)$，并注意
$N_p^\sigma(S W_f^s(y))\cong N_p^\sigma W_f^s(y)$，得到
\[
\begin{aligned}
 T_pM
 &=o_\sigma(x,y;f)T_p^\sigma(S W_f^u(x))\oplus T_p^\sigma W_f^s(y)\\
 &=o_\sigma(x,y;f)(-1)^kT_p^\sigma W_f^u(x)
 \oplus T_p^\sigma(S W_f^s(y))\\
 &=o_\sigma(x,y;f)(-1)^k(-1)^{(k+1)(n-k-1)}
 T_p^\sigma(S W_f^s(y))\oplus T_p^\sigma W_f^u(x).
\end{aligned}
\]
故由~\eqref{eq:19-stable-minus-f-unstable-f}，
\[
\begin{aligned}
 &(-1)^{k+1+n-k-1}o_{\widetilde\sigma}(y,x;-f)
 T_p^{\widetilde\sigma}(S W_{-f}^u(y))\\
 &\qquad=(-1)^{k+(k+1)(n-k-1)}o_\sigma(x,y;f)
 T_p^\sigma(S W_f^s(y)).
\end{aligned}
\]
最后注意
\[
 \{\nabla f(p)\}\oplus T_p^{\widetilde\sigma}(S W_{-f}^u(y))
 =T_p^{\widetilde\sigma}W_{-f}^u(y)
\]
以及
\[
 \{-\nabla f(p)\}\oplus T_p^\sigma(S W_f^s(y))=T_p^\sigma W_f^s(y).
\]
把上述关系移到 $y$，得到
\begin{equation}
\begin{aligned}
 &(-1)^{k+1+n-k-1}o_{\widetilde\sigma}(y,x;-f)
 T_y^{\widetilde\sigma}W_{-f}^u(y)\\
 &\qquad=(-1)^{k+(k+1)(n-k-1)+1}o_\sigma(x,y;f)
 T_y^\sigma W_f^s(y).
\end{aligned}
\label{eq:19-orientation-transfer-to-critical-point}
\end{equation}
又有
\[
 N_y^{\widetilde\sigma}W_{-f}^u(y)
 \oplus T_y^{\widetilde\sigma}W_{-f}^u(y)=T_yM
\]
与
\[
 N_y^\sigma W_f^s(y)\oplus T_y^\sigma W_f^s(y)=T_yM.
\]
因此只需比较 $N_y^{\widetilde\sigma}W_{-f}^u(y)$ 与 $N_y^\sigma W_f^s(y)$。
为此直和 $T_y^{\widetilde\sigma}W_{-f}^u(y)$ 并计算
\footnote{译者注：原书这段计算全在临界点 $y$ 的切空间中进行，末项却印作
$T_pM$；译文改为 $T_yM$。}
\[
\begin{aligned}
 N_y^\sigma W_f^s(y)\oplus T_y^{\widetilde\sigma}W_{-f}^u(y)
 &=N_y^\sigma W_f^s(y)\oplus N_y^{\widetilde\sigma}W_{-f}^s(y)\\
 &=T_y^\sigma W_f^u(y)\oplus N_y^\sigma W_f^u(y)\\
 &=(-1)^{k(n-k)}T_yM,
\end{aligned}
\]
其中倒数第二个等号来自 $\widetilde\sigma$ 的定义。这证明
\[
 T_y^{\widetilde\sigma}W_{-f}^u(y)
 =(-1)^{k(n-k)}T_y^\sigma W_f^s(y).
\]
代入~\eqref{eq:19-orientation-transfer-to-critical-point}，得到
\[
 (-1)^{k+1+n-k-1}o_{\widetilde\sigma}(y,x;-f)
 =(-1)^{k+(k+1)(n-k-1)+1+k(n-k)}o_\sigma(x,y;f).
\]
化简指数，便有
\[
 o_{\widetilde\sigma}(y,x;-f)=(-1)^ko_\sigma(x,y;f),
\]
证明完毕。
\end{proof}

由此构造出如下典范同构。

\begin{corollary}[时间反演]
\label{cor:morse-time-reversal-duality}
对任意不必可定向的流形 $M$ 及所有 $k\in\mathbb N$，
\begin{equation}
 HM_{n-k}(-f;\widetilde\sigma)\cong HM^k(f;\sigma).
\label{eq:19-morse-time-reversal-duality}
\end{equation}
\end{corollary}

\subsection{\personMorse{}同调中的\personPoincare{}对偶}
\label{subsec:poincare-duality-morse-homology}

现设 $M$ 可定向，并给定定向 $o_M$。回忆，对有向子流形 $S\subset M$，其法丛在
$p\in S$ 处的定向由定向向量空间关系
\begin{equation}
 N_pS\oplus T_pS=T_pM
\label{eq:19-normal-tangent-ambient-orientation}
\end{equation}
确定。特别地，若按照第 15.6.1 节的方案给 $T_yW_y^u(f)$ 赋予定向，则关系
\[
 N_yW_y^u(f)\oplus T_y^\sigma W_y^u(f)=T_yM
\]
典范地给法空间 $N_yW_y^u(f)$ 本身赋予定向；\emph{没有环境定向时则不能如此}。在当前
情形，由于 $N_yW_f^s(y)=T_yW_f^u(y)$ 已按 $\sigma$ 的定向方案定向，\emph{稳定}流形
$W_f^s(y)$ 也由标准关系
\begin{equation}
 N_y^\sigma W_f^s(y)\oplus T_yW_f^s(y)=T_yM
\label{eq:19-orient-stable-manifold}
\end{equation}
获得自然定向。

把相应的定向向量空间 $T_yW_f^s(y)$ 记作 $T_y^{\check\sigma}W_f^s(y)$。再要求
\begin{equation}
 N^{\check\sigma}W_f^u(y)=T_y^{\check\sigma}W_f^s(y),
\label{eq:19-coorient-unstable-by-stable}
\end{equation}
从而给\emph{不稳定流形}配上余定向。因此，若 $M$ 可定向，则关系
\eqref{eq:19-orient-stable-manifold}--\eqref{eq:19-coorient-unstable-by-stable}
还让 $\sigma(f)$ 产生另一个系统 $\check\sigma(f)$；它给 $f$ 本身而非 $-f$ 的稳定
流形赋予定向，并给不稳定流形赋予余定向。现在比较给定\personMorse{}函数 $f$ 上的
$\sigma(f)$ 与 $\check\sigma(f)$。

系统 $\widetilde\sigma$ 给 $f$ 的稳定流形（或 $-f$ 的不稳定流形）赋予定向。要求
定向向量空间关系
\[
 N_y^{\widetilde\sigma}W_{-f}^u(y)=T_y^{\check\sigma}W_{-f}^s(y),
\]
便在每个 $y$ 处诱导 $N_y^{\check\sigma}W_{-f}^u(y)$ 的一个定向。

\begin{proposition}
\label{prop:compare-two-morse-orientation-systems}
设 $M$ 可定向且 $n=\dim M$，则
\[
 o_{\check\sigma}(x,y;f)=(-1)^no_\sigma(x,y;f).
\]
特别地，符号只依赖于 $\dim M$，而不依赖于同一指标的各个临界点。
\end{proposition}

\begin{proof}
设 $\sigma$ 为上述相容系统。由可定向性和 $\check\sigma$ 的定义，
\[
 T_x^{\check\sigma}W_f^s(x)\oplus T_x^\sigma W_f^u(x)=T_xM,
 \qquad T_x^{\check\sigma}W_f^s(x)=N_x^{\check\sigma}W_f^u(x).
\]
符号 $o_{\check\sigma}(x,y;f)$ 由关系
\[
 N_p^{\check\sigma}(S W_f^u(x))
 =o_{\check\sigma}(x,y;f)T_p^{\check\sigma}(S W_f^s(y))
\]
确定。右侧直和 $T_p^{\check\sigma}(S W_f^u(x))$，得到
\begin{equation}
 T_p^{\check\sigma}M^c
 =o_{\check\sigma}(x,y;f)T_p^{\check\sigma}(S W_f^s(y))
 \oplus T_p^{\check\sigma}(S W_f^u(x)).
\label{eq:19-check-orientation-transverse-slice}
\end{equation}

另一方面，$o_\sigma(x,y;f)$ 由
\[
 N_p^\sigma(S W_f^s(y))=o_\sigma(x,y;f)T_p^\sigma(S W_f^u(x))
\]
确定。右侧直和 $T_p^\sigma S W_f^s(y)$，得到
\begin{equation}
 T_p^\sigma M^c=o_\sigma(x,y;f)T_p^\sigma(S W_f^u(x))
 \oplus T_p^\sigma(S W_f^s(y)).
\label{eq:19-sigma-orientation-transverse-slice}
\end{equation}

由 $\check\sigma$ 的定义得到如下引理；其证明只是直接计算，留给读者。

\begin{lemma}
\label{lem:check-sigma-orientation-comparison}
若 $\mu_f(y)=k$ 且 $\mu_f(x)=k+1$，则
\footnote{译者注：原书第一式右端仍印作
$T^{\check\sigma}W_f^s(y)$，使等式成为同一定向等于自身乘一个未必为正的符号；
随后代入式 (19.7.37) 时实际使用的是与 $T^\sigma W_f^s(y)$ 的比较。译文据该计算
改正右端上标。}
\[
 T^{\check\sigma}W_f^s(y)=(-1)^{k(n-k)}T^\sigma W_f^s(y),
 \qquad
 T^{\check\sigma}W_f^u(x)=(-1)^{n-k-1}T_p^\sigma W_f^u(x).
\]
\end{lemma}

引理第二式也可写成
\[
 \{\nabla f(p)\}\oplus T^{\check\sigma}S W_f^u(x)
 =(-1)^{n-k-1}\{-\nabla f(p)\}\oplus T_p^\sigma S W_f^u(x),
\]
因而
\[
 T^{\check\sigma}S W_f^u(x)=(-1)^{n-k}T_p^\sigma S W_f^u(x).
\]
代入~\eqref{eq:19-check-orientation-transverse-slice}，得到
\begin{equation}
\begin{aligned}
 T_pM
 &=\{\nabla f(p)\}\oplus T_p^{\check\sigma}M^c\\
 &=o_{\check\sigma}(x,y;f)\{\nabla f(p)\}
 \oplus T_p^{\check\sigma}(S W_f^s(y))
 \oplus T_p^{\check\sigma}(S W_f^u(x))\\
 &=o_{\check\sigma}(x,y;f)T_p^{\check\sigma}W_f^s(y)
 \oplus T_p^{\check\sigma}(S W_f^u(x))\\
 &=(-1)^{k(n-k)+n-k}T_p^\sigma W_f^s(y)
 \oplus T_p^\sigma(S W_f^u(x))\\
 &=(-1)^{k(n-k)+n-k+k(n-k)}T_p^\sigma(S W_f^u(x))
 \oplus T_p^\sigma W_f^s(y).
\end{aligned}
\label{eq:19-check-orientation-total-sign}
\end{equation}
另一方面，式~\eqref{eq:19-sigma-orientation-transverse-slice} 给出
\begin{equation}
\begin{aligned}
 T_pM
 &=\{-\nabla f\}\oplus T_p^\sigma M^c\\
 &=o_\sigma(x,y;f)\{-\nabla f\}\oplus T_p^\sigma(S W_f^u(x))
 \oplus T_p^\sigma(S W_f^s(y))\\
 &=o_\sigma(x,y;f)(-1)^kT_p^\sigma(S W_f^u(x))
 \oplus\{-\nabla f\}\oplus T_p^\sigma(S W_f^s(y))\\
 &=o_\sigma(x,y;f)(-1)^kT_p^\sigma(S W_f^u(x))
 \oplus T_p^\sigma W_f^s(y).
\end{aligned}
\label{eq:19-sigma-orientation-total-sign}
\end{equation}
比较~\eqref{eq:19-check-orientation-total-sign} 与
\eqref{eq:19-sigma-orientation-total-sign} 并化简指数，得到
\[
 o_{\check\sigma}(x,y;f)=(-1)^no_\sigma(x,y;f),
\]
证明完毕。
\end{proof}

特别地，对任意 $f:M\to\mathbb R$，
\begin{equation}
 \left.\partial_f^{\check\sigma}\right|_{C_k(f)}
 =(-1)^n\left.\partial_f^\sigma\right|_{C_k(f)},
\label{eq:19-boundary-orientation-system-sign}
\end{equation}
从而
\begin{equation}
 H_{n-k}^{\check\sigma}(f;\mathbb Z)
 \cong H_{n-k}^\sigma(f;\mathbb Z).
\label{eq:19-morse-homology-orientation-systems-isomorphic}
\end{equation}
显然，相应上同调也有同一同构。结合
\eqref{eq:19-morse-time-reversal-duality} 与
\eqref{eq:19-morse-homology-orientation-systems-isomorphic}，得到
\begin{equation}
 H^k(f;\sigma)\cong H_{n-k}^\sigma(f;\mathbb Z).
\label{eq:19-poincare-duality-morse}
\end{equation}
在同构
\begin{equation}
 H_*^\sigma(f;M)\cong H_*^{\mathrm{sing}}(M;\mathbb Z)
\label{eq:19-morse-singular-homology-isomorphism}
\end{equation}
下，这恰好与经典\personPoincare{}对偶同构 $a\mapsto a\cap[M]$ 一致。（在经典\personMorse{}理论
情形下，含符号的证明见 (Mil65)；它本质上与上述证明相同。）我们强调，时间反演
同构~\eqref{eq:19-morse-time-reversal-duality} 并不总要求 $M$ 有定向，而\personPoincare{}对偶同构~\eqref{eq:19-poincare-duality-morse} 要求 $M$ 可定向。

事实上，可以用~\eqref{eq:19-poincare-duality-morse} 和
\eqref{eq:19-morse-singular-homology-isomorphism} 找到满足
\begin{equation}
 H_*^\sigma(f;M)\not\cong H_*^{\check\sigma}(f;M)
\label{eq:19-two-morse-homologies-differ}
\end{equation}
的 $M$。（以下讨论取自 (Oh99)。）
\footnote{译者注：原书式 (19.7.45) 比较的是 $\sigma$ 与
$\check\sigma$，但紧接着的例 19.7.8 全程计算并在结尾比较的是 $\sigma$ 与
$\widetilde\sigma$；两种记号在本节定义不同。信息不足以判定作者意欲改哪一处，故
两处均照录并指出不一致。}

\begin{example}
\label{ex:even-real-projective-space-orientations}
考虑偶数维实射影空间 $\mathbb{RP}^{2m}$。它不可定向，并且
\footnote{译者注：原书此处断言
$H_{\mathrm{sing}}^{2m}(\mathbb{RP}^{2m};\mathbb Z)=0$，并据此在下文得到
$H_0^{\widetilde\sigma}=0$。但当 $m\geq1$ 时，泛系数定理给出
$H^{2m}(\mathbb{RP}^{2m};\mathbb Z)\cong\mathbb Z/2$；因此原书的这个上同调值及其
下游结论有误。为保留本例原有推演，正文照录并在此指出。}
\[
 H_0^{\mathrm{sing}}(\mathbb{RP}^{2m};\mathbb Z)\cong\mathbb Z,
 \qquad H_{2m}^{\mathrm{sing}}(\mathbb{RP}^{2m};\mathbb Z)=0,
\]
\[
H_{\mathrm{sing}}^0(\mathbb{RP}^{2m};\mathbb Z)\cong\mathbb Z,
 \qquad H_{\mathrm{sing}}^{2m}(\mathbb{RP}^{2m};\mathbb Z)=0.
\]
因此，由~\eqref{eq:19-morse-singular-homology-isomorphism}，
\[
 H_0^\sigma(f;\mathbb{RP}^{2m})\cong\mathbb Z,
 \qquad H_{2m}^\sigma(f;\mathbb{RP}^{2m})=0.
\]
另一方面，式~\eqref{eq:19-morse-time-reversal-duality} 蕴含
\[
 H_0^{\widetilde\sigma}(-f;\mathbb{RP}^{2m})
 \cong H_{\mathrm{sing}}^{2m}(\mathbb{RP}^{2m};\mathbb Z)=0,
\]
\[
 H_{2m}^{\widetilde\sigma}(-f;\mathbb{RP}^{2m})
 \cong H_{\mathrm{sing}}^0(\mathbb{RP}^{2m};\mathbb Z)=\mathbb Z.
\]
只要固定相容定向系统，延拓同构便可把这个同构转移到 $f$ 本身，故
\[
 H_0^{\widetilde\sigma}(f;\mathbb{RP}^{2m})
 \cong H_{\mathrm{sing}}^{2m}(\mathbb{RP}^{2m};\mathbb Z)=0,
\]
\[
 H_{2m}^{\widetilde\sigma}(f;\mathbb{RP}^{2m})
 \cong H_{\mathrm{sing}}^0(\mathbb{RP}^{2m};\mathbb Z)=\mathbb Z.
\]
这个例子表明，两个\personMorse{}同调 $H^\sigma(f;M)$ 与
$H^{\widetilde\sigma}(f;M)$ 的确可能不同。
\end{example}

\subsection{\personFloer{}复形中的对偶性}
\label{subsec:duality-floer-complex}

受上一小节\personMorse{}复形讨论的启发，我们考察\personFloer{}理论中类似的时间反演映射的
作用。两者一个突出的差别是：这里需要研究\emph{无限维}情形中的对偶性。

从自由环路空间 $\mathcal L_0(M)$ 及其\personNovikov{}覆盖
$\widetilde{\mathcal L}_0(M)$ 上的 $\mathbb Z_2$-对称开始：
\begin{equation}
 [z,w]\longmapsto[\widetilde z,\widetilde w],
\label{eq:19-loop-time-reversal-involution}
\end{equation}
其中 $\widetilde z(t)=z(1-t)$，而
$\widetilde w(x)=w(\bar x)$，$\bar x$ 是 $x\in D^2\subset\mathbb C$ 的复共轭。
还考虑 $H,J$ 及相容定向系统上的操作
\[
 H\longmapsto\widetilde H,\qquad
 J\longmapsto\widetilde J,\qquad
 \sigma\longmapsto\widetilde\sigma,
\]
其中
\[
 \widetilde H(t,x)=-H(1-t,x),\qquad \widetilde J_t=J_{1-t}.
\]
由\personFloer{}方程容易验证，这个对称在相应\personFloer{}方程的解集之间诱导一一对应，并给出
矩阵系数的恒等式
\[
 n_{(H,J)}^\sigma([z,w],[z',w'])
 =n_{(\widetilde H,\widetilde J)}^{\widetilde\sigma}
 ([\widetilde z',\widetilde w'],[\widetilde z,\widetilde w]).
\]
这诱导单射链同态
\[
 CF^{\widetilde\sigma}(\widetilde H;\mathbb Q)
 \lhook\joinrel\longrightarrow
 \operatorname{Hom}(CF^\sigma(H);\mathbb Q).
\]

在生成元上令
\[
 \left\langle[\widetilde z,\widetilde w],[z',w']\right\rangle
 =\begin{cases}
 q^{-\left(\int\widetilde w^*\omega+
 \int(w')^*\omega\right)},&z=z',\\
 0,&z\neq z',
 \end{cases}
\]
并双线性延拓，从而定义自然非退化配对。这里指数看似不依赖哈密顿量 $H$，但这是
误导性的；当 $z'=z$ 时，对任意 $H$ 都有普适公式
\begin{equation}
 \mathcal A_{\widetilde H}([\widetilde z,\widetilde w])
 +\mathcal A_H([z',w'])
 =-\int\widetilde w^*\omega-\int(w')^*\omega.
\label{eq:19-time-reversed-action-sum}
\end{equation}
这是因为哈密顿项在如下变量替换中抵消：
\[
\begin{aligned}
 \int_0^1\widetilde H(t,\widetilde z(t))\,dt
 &=\int_0^1-H(1-t,z(1-t))\,dt\\
 &=-\int_0^1H(t,z(t))\,dt
 =-\int_0^1H(t,z'(t))\,dt.
\end{aligned}
\]

有显式公式
\[
 \langle\widetilde\alpha,\beta\rangle
 =\sum_{\substack{A\\c_1(A)=0}}
 \left(\sum_{[z,w]\in\operatorname{Crit}\mathcal A_H}
 a_{[\widetilde z,\widetilde w]}
 b_{[z,w\#(-A)]}\right)q^A,
\]
其中
\[
 \widetilde\alpha
 =\sum_{[\widetilde z,\widetilde w]\in
 \operatorname{Crit}\mathcal A_{\widetilde H}}
 a_{[\widetilde z,\widetilde w]}[\widetilde z,\widetilde w],
 \qquad
 \beta=\sum_{[z,w]\in\operatorname{Crit}\mathcal A_H}b_{[z,w]}[z,w].
\]
此外，\personConley{}--\personZehnder{}指标满足
\[
 \mu_H([z,w])=2n-\mu_{\widetilde H}([\widetilde z,\widetilde w]),
\]
因此该配对诱导单射同态
\[
 CF_k^{\widetilde\sigma}(\widetilde H)
 \lhook\joinrel\longrightarrow
 \operatorname{Hom}_{\Lambda_\omega^{(0)}}
 \bigl(CF_{2n-k}^{\sigma}(H);\Lambda_\omega^{(0)}\bigr),
\]
其中
\[
 \Lambda_\omega^{(0)}=
 \left\{\sum_Aa_Aq^A=\sum_Aa_AT^{\omega(A)}\ \middle|\ c_1(A)=0\right\}.
\]
注意，$\Lambda_\omega^{(0)}$ 是域，而
$CF_k^{\widetilde\sigma}(\widetilde H)$ 与 $CF_{2n-k}^\sigma(H)$ 都是其上的有限维
向量空间，且生成元数相同。因此由泛系数定理，此同态诱导同构
\begin{equation}
 HF_k(\widetilde H,\widetilde J;\widetilde\sigma)
 \cong\operatorname{Hom}_{\Lambda_\omega^{(0)}}
 \bigl(HF_{2n-k}(H,J;\sigma),\Lambda_\omega^{(0)}\bigr).
\label{eq:19-floer-time-reversal-dual-isomorphism}
\end{equation}
这里回忆，\personFloer{}模空间上的相容定向系统 $\widetilde\sigma$ 由
$\mathcal L_0(M)$ 上的映射~\eqref{eq:19-loop-time-reversal-involution} 典范地诱导。

总结上述讨论如下。

\begin{theorem}
\label{thm:strongly-nondegenerate-floer-pairing}
设 $\Lambda_\omega^{(0)}$ 如上，则存在强非退化配对
\begin{equation}
 L=\langle\,\cdot,\cdot\,\rangle:
 HF_k(\widetilde H,\widetilde J;\widetilde\sigma)
 \times HF_{2n-k}(H,J;\sigma)
 \longrightarrow\Lambda_\omega^{(0)}.
\label{eq:19-strong-floer-pairing}
\end{equation}
\end{theorem}

接着定义投影
\begin{equation}
 \pi_0:\Lambda_\omega^{(0)}\longrightarrow\mathbb Q;\qquad
 \pi_0(\lambda)=\lambda_0,
\label{eq:19-zero-novikov-coefficient-projection}
\end{equation}
其中 $\lambda_0\in\mathbb Q$ 是
$\lambda=\sum_{A\in\Gamma}\lambda_Aq^A$ 中 $A=0$ 的系数。沿用 (Os03) 的记号
\footnote{译者注：原书在投影定义后的级数中把系数下标排成
$\lambda_{A\in\Gamma}$；译文按求和指标还原为 $\lambda_A$。}
，令
\begin{equation}
 L'=\pi_0\circ L:
 HF_k(\widetilde H,\widetilde J;\widetilde\sigma)
 \times HF_{2n-k}(H,J;\sigma)\longrightarrow\mathbb Q.
\label{eq:19-rational-floer-pairing}
\end{equation}

目前尚不明显 $L'$ 是否非退化。第 20.4 节将依照\personEntov{}与\personPolterovich{}
(EnP03)、(Os03) 的方案，对有理辛流形证明 $L'$ 的非退化性；该方案把此配对与量子
上同调环中的\personFrobenius{}配对联系起来。无理情形的非退化性需要使用非平凡的
非阿基米德分析论证，由\personUsher{} (Ush10b) 证明。
